\documentclass[11pt,letterpaper]{article}
\usepackage[T1]{fontenc}
\usepackage{lmodern}
\usepackage[margin=1.05in]{geometry}
\usepackage{amsmath,amssymb,amsthm,booktabs,microtype,xurl}
\usepackage[hidelinks,pdftitle={Nonhalting and PA-independence for a 196-state Turing machine},
  pdfauthor={James Chang}]{hyperref}

\newtheorem{theorem}{Theorem}[section]
\newtheorem{lemma}[theorem]{Lemma}
\newcommand{\PA}{\mathrm{PA}}
\newcommand{\N}{\mathbb{N}}
\newcommand{\Con}{\operatorname{Con}}
\newcommand{\Prov}{\operatorname{Prov}}
\newcommand{\ax}{\mathsf{AX}}
\newcommand{\mpc}{\mathsf{MP}}
\newcommand{\shift}{\mathsf{S}}
\newcommand{\Ic}{\mathsf{I}}
\newcommand{\Mc}{\mathsf{M}}
\newcommand{\Uc}{\mathsf{U}}
\newcommand{\Ec}{\mathsf{E}}
\newcommand{\Nc}{\mathsf{N}}
\newcommand{\Rc}{\mathsf{R}}
\newcommand{\file}[1]{\texttt{\nolinkurl{#1}}}
\setlength{\emergencystretch}{2em}

\title{Nonhalting and PA-independence\\for a 196-state Turing machine}
\author{James Chang}
\date{}

\begin{document}
\maketitle

\begin{abstract}
We give a two-symbol Turing machine with 196 working states that
searches for a contradiction in ordinary first-order Peano arithmetic.
It runs forever on a blank tape. Peano arithmetic proves neither its
exact arithmetic halting sentence nor its negation. The machine's
axiom command reverses the two arguments in the quantifier code and
uses a guarded equality-vacuity clause in place of a primitive
transitivity head. In Lean~4 we verify the calculus, the simulation
of the register program, and both unprovability results for the
literal table.
\end{abstract}

\section{The result}\label{sec:result}

Let $M$ be the machine in Section~\ref{sec:machine}, and let $M_{209}$
be its predecessor~\cite{pa209}. Write $\PA$ for ordinary first-order
Peano arithmetic. The arithmetic halting sentence $H$ uses the rows
of $M$ in the trace formula of~\cite[\S5]{pa209}.

\begin{theorem}\label{thm:main}
If $\PA$ is consistent, the machine $M$ does not halt from the blank
tape and $\PA\nvdash\neg H$. If $\PA$ is sound for the standard model,
also $\PA\nvdash H$.
\end{theorem}

The change to the register program is confined to the body of one
command (Lemma~\ref{lem:program}). We prove soundness for the revised
calculus in Section~\ref{sec:sound} and show how the compiler
of~\cite{pa209} sends PA refutations to proofs in that calculus in
Section~\ref{sec:complete}. Sections~\ref{sec:layout}
and~\ref{sec:simulation} give the tape layout and its simulation proof.

Lean uses Foundation's standard model of ordinary PA to establish
consistency and soundness in the metatheory. This yields closed Lean
theorems for blank-tape nonhalting and both unprovability directions.

\section{The machine}\label{sec:machine}

The tape has alphabet $\{0,1\}$ and is infinite in both directions.
It starts with all cells zero. State $0$ is the start state,
$0,\ldots,195$ are the working states, and $196$ is the halt label. In state $q$, a row
$(q,b,w,D,q')$ reads $b$, writes $w$, moves one cell in direction $D$,
and enters $q'$. The table has one row for each pair in
$\{0,\ldots,195\}\times\{0,1\}$. The file \file{data/pa196.tm} contains
all 392 rows and has SHA-256 digest
\begin{quote}\small
\nolinkurl{e8bbccd4cfa325251ca95066db116b5aa766e404b0971c33c02772be1c3617a8}.
\end{quote}
The first row is $(0,0,1,R,1)$, and $(63,1,1,R,196)$ is the only row
that enters the halt label.

As in $M_{209}$, a tape engine interprets a register program stored
as a binary tree. A \emph{frame} records a marker, the path to the
current site, a gap, and the register file. A block $0\,1^{v+1}$
stores value $v$~\cite[\S3]{pa209}. The engine's 36 states have
the same rows as in $M_{209}$ after renaming. Of the remaining states,
13 route leaf instructions to the engine and 147 read the 159
distinct nodes. State $75$ reads the root.

\section{The register program}\label{sec:steps}

The register program of $M_{209}$ enumerates words over three commands
on a store $(t,w,d,h)$. It halts when Shift reads formula code
$1$~\cite[\S4]{pa209}. We keep Shift and $\mpc$, change $\ax(c,x,y)$
in three steps, and then change how the program stores that command.

\subsection{Step 1. The quantifier code}

Let $p(a,b)=\binom{a+b+1}{2}+b$. The four constructors in $M_{209}$
use $4p(a,b)+r$, where $r=0,1,2,3$ denotes implication, membership,
universal quantification and equality, respectively. The constructors
in $M$ are
\[
 \Ic(a,b)=4p(a,b),\quad \Mc(a,b)=4p(a,b)+1,\quad
 \Uc(a,B)=4p(B,a)+2,\quad \Ec(a,b)=4p(a,b)+3 .
\]
Only the quantifier code reverses the pairing arguments. As before,
code $0$ is truth and $1=\Mc(0,0)$ is the contradiction. Write
$\Nc(B)=\Ic(B,1)$, $B\wedge C=\Nc(\Ic(B,\Nc(C)))$,
$B\leftrightarrow C$ for $\Ic(B,C)\wedge\Ic(C,B)$, and
$\exists x_a\,B$ for $\Nc(\Uc(a,\Nc(B)))$. An arrow in the formulas
below denotes $\Ic$.

\subsection{Steps 2 and 3. The axiom heads}

In $M_{209}$, $\ax$ pushes the saved top, a generalization, the
restricted propositional head $\rho_8$ and seven further heads
$\rho_1,\ldots,\rho_7$. It leaves a guarded gate in $t$.
For $\ax(c,x,y)$ in $M$, put $D=d+x$, $H=h+y$ and $A=c$. The command
pushes the saved top $t$, the generalization $\Uc(d,h)$, the head
\[
 \rho_8(D,A,H,d)=\bigl((H\to d)\to D\bigr)\to\bigl((D\to H)\to(A\to H)\bigr),
\]
and then the six formulas of Table~\ref{tab:heads}, in that order.

\begin{table}[ht]
\centering\small
\begin{tabular}{@{}cll@{}}
\toprule
Push & Formula & Head of~\cite{pa209} \\
\midrule
1 & $(\exists x_D\,\Uc(D,H))\to H$ & $\rho_6$ \\
2 & $\Uc(D,\Uc(A,H))\to\Uc(A,\Uc(D,H))$ & $\rho_4$ \\
3 & $\Mc(H,D)\to(\Ec(H,A)\to\Mc(A,D))$ & $\rho_5$, antecedents exchanged \\
4 & $\Ec(H,D)\to(\Mc(A,H)\to\Mc(A,D))$ & $\rho_3$ \\
5 & $\exists x_D\,\Ec(H,D)$ & $\rho_1$, equation mirrored \\
6 & $\Uc(D,A)\to(\Uc(D,A\to H)\to\Uc(D,H))$ & $\rho_2$, antecedents exchanged \\
\bottomrule
\end{tabular}
\caption{The six heads pushed by $\ax$ after the saved top,
$\Uc(d,h)$ and $\rho_8$.}\label{tab:heads}
\end{table}

The command leaves $0$ in $t$ if $A=H$ or $A=D$, and $F(H,D,A)$
otherwise. For labels $a,b$ and a code $c$, the gate $F$ is defined by
\begin{align*}
 \Delta_c(a,b)&=\Ec(a,b)\to\Uc(c,\Ec(a,b)), &
 V&=\Mc(a,b)\to\Uc(c,\Mc(a,b)),\\
 P&=c\to\Mc(c,b), &
 S&=\Uc\bigl(b,\ (\Mc(a,b)\to P)\wedge(\Mc(c,b)\to c)\bigr),\\
 Q&=\Uc(c,\,S\to\Nc(P)), &
 F(a,b,c)&=\Nc\bigl(\Delta_c(a,b)\to(V\to Q)\bigr).
\end{align*}
The command finishes by setting $d=h=0$. It pushes nine items in all,
including the saved top. In $P$, $S$ and $Q$, the number $c$ is both
the label of $x_c$ and a formula code.

Under the guard $c\ne a,b$, the gate $G(a,b,c)=\Nc(V\to Q)$
of~\cite{pa209} gives membership vacuity $V$ and pointed separation
$\Nc(Q)$. Step~2 removes the seventh head
\[
 \rho_7(a,b,c)=\Ec(c,a)\to\bigl(\Ec(c,b)\to\Uc(\Mc(a,b),\Ec(a,b))\bigr)
\]
and adds equality vacuity $\Delta_c(a,b)$ to the gate under the same
guard. Step~3 rearranges the heads as shown in the third column of
Table~\ref{tab:heads} and exchanges the second and third arguments of
$\rho_8$.

The body of $\ax$ is a straight-line sequence of 201 items. Its macro
calls come from~\cite{pa209}. They push the accumulator, push it and
load one of the three registers, pop and form $4p$, read the two
parameters, or check the two parts of the guard. The other items are
single increments of the accumulator and the closing jump. The full
program has $3{,}812$ sites, compared with $3{,}948$ for $M_{209}$.

\begin{lemma}[Program]\label{lem:program}
Outside the body of $\ax$ the register program $R$ of $M$ agrees with
that of $M_{209}$ after translation of addresses. The body of $\ax$ is a concatenation of blocks, each
identical to a macro block of $M_{209}$ and left only to the next
block. It computes the saved top, $\Uc(d,h)$, $\rho_8(D,A,H,d)$, the
six formulas of Table~\ref{tab:heads} and the gate $F(H,D,A)$, in the
codes of Step~1.
\end{lemma}

\begin{proof}
After flattening the program and translating the addresses, the 169
sites outside the body agree with $M_{209}$. The 201 blocks use the
same macros and internal jump offsets. The closing jump goes to the
corresponding successor. The macro contracts give the store effect of
each block. Symbolic execution gives the nine stack entries and the
new top as formula trees over $4p(\cdot,\cdot)+r$, with the argument
order of Step~1. Lean checks the source program in
\file{PA196/Program.lean} and the effect of the command in
\file{PA196/Source/AXEffect.lean}.
\end{proof}

\section{Soundness}\label{sec:sound}

\begin{lemma}\label{lem:sound}
If $\PA$ is consistent, $R$ does not halt from the clean store.
\end{lemma}

\begin{proof}
Use the interpretation of membership and equality
from~\cite[\S4.2]{pa209}. In a model of $\PA$, $\Mc(a,b)$ says that
bit $x_b$ of $x_a$ is set. Each formula in Table~\ref{tab:heads} is
either a head of~\cite{pa209}, possibly with two antecedents exchanged,
or the witness formula $\exists x_a(x_b=x_a)$. Each is true under every
assignment, as is the propositional tautology $\rho_8$.

When $c\ne a,b$, $G(a,b,c)$ is true by~\cite{pa209}. The atom
$\Ec(a,b)$ contains no $x_c$, so $\Delta_c(a,b)$ is also true. Thus
$F(a,b,c)$ is true, since it is propositionally $G\wedge\Delta_c$.
If the guard fails, the gate returns truth. Modus ponens and
generalization preserve truth under every assignment, and code $1$
is false. We use the invariant, halt test and scheduler
of~\cite[\S4.2]{pa209}.

The encoder and decoder exchange the quantifier arguments together.
For a nonzero code $4p(u,v)+r$, both components are smaller than the
code in either order. The decoder still recurses on smaller codes, and
every label in the formula with code $c$ is below $c$. Thus $x_c$ is
fresh in the gate.
\end{proof}

\section{Completeness}\label{sec:complete}

A formula is \emph{derivable} if modus ponens and generalization give
a proof from instances of the six heads, $\rho_8$ with $u\le a$, and
the gate with $c\ne a,b$. The proofs below use fixed finite templates
in the labels and codes. Their case distinctions compare numbers
for equality. PA can carry out these arguments on coded proofs by the
method of~\cite[\S6.1]{pa209}.

\subsection{Uses of the seventh head}

In the Lean development of~\cite{pa209}, only
\file{ActualCertificateNeighborhoodPrimitive.lean} uses $\rho_7$
directly. It derives reflexivity and symmetry of primitive equality
and vacuity of $\Ec(a,b)$ under the label $\Mc(a,b)$.
\file{ActualCertificateNeighborhoodCongruence.lean} uses symmetry
to reverse the membership rewrites $\rho_3,\rho_5$.
\file{ActualCertificateNeighborhoodWitness.lean} uses symmetry and
vacuity to prove $\Ec(a,b)\to\Rc(a,b,z)$ for $z\ne a,b$, where
\[
 \Rc(a,b,z)=\Uc\bigl(z,\Mc(a,z)\leftrightarrow\Mc(b,z)\bigr)\wedge
            \Uc\bigl(z,\Mc(z,a)\leftrightarrow\Mc(z,b)\bigr)
\]
is the equality relation of the interpretation. The same module
then gets $\exists x_a\,\Rc(a,c,z)$ from $\rho_1$. The rest of the
proof uses primitive equality through the following three facts.
\begin{description}
\item[(R)] For all labels $a,b,d$, both
  $\Ec(a,b)\to(\Mc(d,a)\leftrightarrow\Mc(d,b))$ and
  $\Ec(a,b)\to(\Mc(a,d)\leftrightarrow\Mc(b,d))$.
\item[(N)] $\Ec(a,b)\to\Rc(a,b,z)$ for $z\ne a,b$.
\item[(W)] An existential witness for equality at every pair of labels.
\end{description}
We derive (R), (N) and (W) directly. They suffice for the interpretation
in place of separate laws for symmetry and transitivity of primitive
equality.

\subsection{Proof rules used below}\label{sec:tools}

\begin{enumerate}
\item The classical implicational rules come from $\rho_8$ as
  in~\cite{pa209}. Use code $1$ for falsity. Head~1 with
  $X=\Uc(D,\Nc(\Uc(D,H)))$ gives $\Nc(X)\to H$. Implicational
  reasoning gives $1\to(X\to1)$, so $1\to H$ for every formula $H$.
  Double negation elimination follows.
\item To lift a derivable implication $B\to C$ to
  $\Uc(x,B)\to\Uc(x,C)$, use generalization and head~6.
\item A formula $B$ is vacuous at $r$ if $B\to\Uc(r,B)$ is
  derivable. The gate gives vacuity for both $\Mc(a,b)$ and
  $\Ec(a,b)$ whenever $r\ne a,b$. If $B$ is vacuous at $r$, lift
  $\Nc(\Uc(r,B))\to\Nc(B)$ at $r$ and use head~1 with $D=r$
  to get $\Nc(\Uc(r,\Nc(B)))\to B$. Thus $\Nc(B)$ is vacuous at $r$.
  Vacuity passes to implications by lifting $\Nc(B)\to(B\to C)$ and
  $C\to(B\to C)$. It passes to $\Uc(x,B)$ when $x\ne r$ by lifting
  at $x$ and using head~2. A formula built from atoms is therefore
  vacuous at every label outside its atom indices and bound labels.
\item If $B$ is vacuous at $x$, then $\Uc(x,A\to B)$ and
  $\exists x\,A$ give $B$. If $\Nc(B)$ holds, vacuity gives
  $\Uc(x,\Nc(B))$, which yields $\Uc(x,\Nc(A))$.
\end{enumerate}
Rules 1, 2 and 4 are from~\cite{pa209}. Rule~3 extends the vacuity
argument of~\cite{pa209} to equality atoms.

\subsection{The gate}

\begin{lemma}\label{lem:extract}
$F(a,b,c)\to G(a,b,c)$ and $F(a,b,c)\to\Delta_c(a,b)$ are derivable.
\end{lemma}

\begin{proof}
Put $B=(V\to Q)$, so $F=\Nc(\Delta_c\to B)$ and $G=\Nc(B)$.
Weakening $B$ gives $\Delta_c\to B$, so $F\to\Nc(B)$. The assumption
$\Nc(\Delta_c)$ also gives $\Delta_c\to B$, by explosion. Therefore
$F\to\Nc(\Nc(\Delta_c))$, and double negation elimination gives
$\Delta_c$ under $F$.
\end{proof}

The new gate thus gives the old gate at the same parameters. For a
label $u$ and a formula $\varphi$ with code $c\ne u$, the instance
at $(u,u,c)$ gives separation in the form
\begin{equation}\label{eq:sep}
 \exists x_c\,(S\wedge P),\qquad
 P=\varphi\to\Mc(c,u),\qquad
 S\to\Uc\bigl(u,\ \Mc(c,u)\to\varphi\bigr).
\end{equation}
The variable $x_u$ is bound in $S$ and free in $P$.

\subsection{Reverse rewriting}

Heads 4 and 3 give the direct rewrites
$\Ec(a,b)\to(\Mc(d,a)\to\Mc(d,b))$ and
$\Ec(a,b)\to(\Mc(a,d)\to\Mc(b,d))$ for all labels. Let $\chi(t)$ be
one of the atoms $\Mc(d,t)$, $\Mc(t,d)$ and $\Mc(t,t)$, where $d$ is a
label. The direct rewrites give
\begin{equation}\label{eq:direct}
 \Ec(v,y)\to(\chi(v)\to\chi(y))
\end{equation}
for all labels $v$ and $y$. For $\Mc(t,t)$, rewrite one position
at a time.

\begin{lemma}[Specialization of a negative property]\label{lem:spec}
Let $y$, $v$, $c$ be labels with $v\ne y$ and $c\ne y$, and let
$\chi$ be as above with $d\ne y$. Then
\[
 \Uc\bigl(y,\ \Mc(c,y)\to\Nc(\chi(y))\bigr)\ \to\
 \bigl(\Mc(c,v)\to\Nc(\chi(v))\bigr)
\]
is derivable.
\end{lemma}

\begin{proof}
Write $\psi(t)=\Mc(c,t)\to\Nc(\chi(t))$. Head~4 gives
$\Ec(v,y)\to(\Mc(c,v)\to\Mc(c,y))$, and~\eqref{eq:direct} gives
$\Ec(v,y)\to(\chi(v)\to\chi(y))$. Heads 3 and 4 allow $d=v$
because they have no side condition. These direct rewrites suffice
for $\Ec(v,y)\to(\psi(y)\to\psi(v))$, since $y$ occurs in $\psi(y)$
only in an antecedent and under a negation. Lift at $y$ to obtain
$\Uc(y,\psi(y))\to\Uc(y,\Ec(v,y)\to\psi(v))$. The indices of
$\psi(v)$ lie among $c$, $v$ and $d$, all different from $y$, so
$\psi(v)$ is vacuous at $y$. Apply existential elimination with the
instance $\exists x_y\,\Ec(v,y)$ of head~5 to get $\psi(v)$.
\end{proof}

\begin{lemma}[Reverse rewriting]\label{lem:rev}
For all labels $a,b,d$,
$\Ec(a,b)\to(\Mc(d,b)\to\Mc(d,a))$ and
$\Ec(a,b)\to(\Mc(b,d)\to\Mc(a,d))$ are derivable.
\end{lemma}

\begin{proof}
For $a=b$ both are instances of $B\to(C\to C)$. Let $a\ne b$, and let
$\chi$ be as above with $d\ne a$. We first derive
\begin{equation}\label{eq:chi}
 \Ec(a,b)\to(\chi(b)\to\chi(a)).
\end{equation}
Let $\varphi=\Nc(\chi(a))$. If its code equals $b$, take
$\Nc(\Nc(\Nc(\chi(a))))$ instead, whose code is larger. Let $c$ be the
code of $\varphi$. Then $c$ exceeds $a$ and $d$ and differs from $b$.
Take the separation instance~\eqref{eq:sep} with $u=a$. By
Lemma~\ref{lem:spec} with $y=a$ and $v=b$,
\[
 S\to\bigl(\Mc(c,b)\to\Nc(\chi(b))\bigr),
\]
and head~4 gives $\Ec(a,b)\to(\Mc(c,a)\to\Mc(c,b))$. Hence
$(S\wedge P)\to C$ is derivable for the formula $C$ in~\eqref{eq:chi}.
Suppose $\Ec(a,b)$ and $\chi(b)$ hold and $\chi(a)$ fails. Then
$\varphi$ holds, so $P$ gives $\Mc(c,a)$. The rewrite gives
$\Mc(c,b)$, and $S$ gives $\Nc(\chi(b))$. The indices of $C$ are
$a$, $b$ and $d$, all different from $c$. Vacuity for the equality
atom makes $C$ vacuous at $c$, so existential elimination
with~\eqref{eq:sep} gives $C$.

For $d\ne a$, \eqref{eq:chi} with $\chi(t)=\Mc(d,t)$ and with
$\chi(t)=\Mc(t,d)$ are the two statements. For $d=a$ the statements
are $\Ec(a,b)\to(\Mc(a,b)\to\Mc(a,a))$ and
$\Ec(a,b)\to(\Mc(b,a)\to\Mc(a,a))$. Head~3 gives
$\Ec(a,b)\to(\Mc(a,b)\to\Mc(b,b))$, head~4 gives
$\Ec(a,b)\to(\Mc(b,a)\to\Mc(b,b))$, and~\eqref{eq:chi} with
$\chi(t)=\Mc(t,t)$ gives $\Ec(a,b)\to(\Mc(b,b)\to\Mc(a,a))$.
\end{proof}

In the interpretation, separation gives a set $x_c$ that contains
$x_a$ if $x_a$ lacks $\chi$ and contains only elements that lack
$\chi$. Since $x_b$ belongs to every set that contains $x_a$, if
$x_b$ has $\chi$, then $x_a$ has it too. The calculus specializes a
quantifier with a witness from head~5 and direct rewrites. These
steps supply the specialization without a substitution rule.

\pagebreak[2]
\begin{lemma}[The defined equality]\label{lem:nbhd}
For $z\ne a,b$, $\Ec(a,b)\to\Rc(a,b,z)$ is derivable.
\end{lemma}

\begin{proof}
The direct rewrites and Lemma~\ref{lem:rev} with $d=z$ give
$\Ec(a,b)\to(\Mc(a,z)\leftrightarrow\Mc(b,z))$ and
$\Ec(a,b)\to(\Mc(z,a)\leftrightarrow\Mc(z,b))$. Lift both at $z$ and
compose with $\Ec(a,b)\to\Uc(z,\Ec(a,b))$ from the gate.
\end{proof}

The Lean proofs of reverse rewriting and neighborhood equality are
in \file{PA196/Logic/Reverse.lean} and
\file{PA196/Logic/Neighborhood.lean}. The certificate compiler builds
the same derivations inside arbitrary PA models.

\subsection{Transporting PA refutations}

\begin{lemma}[Completeness]\label{lem:complete}
$\PA$ proves that if $\Prov_{\PA}(\ulcorner\bot\urcorner)$ then $R$
halts from the clean store.
\end{lemma}

\begin{proof}
The direct rewrites and Lemma~\ref{lem:rev} give (R),
Lemma~\ref{lem:nbhd} gives (N), and head~5 gives (W) with the
equation reversed. To specialize a quantifier at $z$ to a label $c$,
the proof in~\cite{pa209} assumes $\Ec(z,c)$ and rewrites atoms in
both directions. With the head~5 witness, the assumption is
$\Ec(c,z)$, and (R) still gives both rewrites. For the witness of the
defined equality, apply existential monotonicity to $\exists x_a\,\Ec(c,a)$ using (N) to get
$\exists x_a\,\Rc(c,a,z)$. The definition of $\Rc$ gives its symmetry,
and hence $\exists x_a\,\Rc(a,c,z)$.

Lemma~\ref{lem:extract} gives the old gate. Heads 1, 2, 4 and
$\rho_8$ give the old heads $\rho_6$, $\rho_4$, $\rho_3$ and
$\rho_8$. Exchanging two antecedents in heads 3 and 6 gives $\rho_5$
and $\rho_2$ by propositional reasoning. With these derivations in
place of the ones that use $\rho_7$, the construction
of~\cite[\S6.1]{pa209} turns a coded $\PA$ refutation into a coded
derivation of code $1$ in the present calculus.

Step~1 preserves the four coding properties the construction uses.
The encoder and decoder are inverse, both components are smaller
than their compound code, the residues modulo $4$ are the same, and
$\mpc$ still computes $\Ic(t,y)=4p(t,y)$. When we choose a formula
code as a label, we use the new constructor. The code of $\varphi$
in Lemma~\ref{lem:rev} and the padded property code for separation
satisfy the same bounds.

The nine items pushed by $\ax$ change the command identities
of~\cite[\S6.2]{pa209}. Writing $\shift$ for Shift, the first is
\[
 \ax(c,x,y)\,\shift^{9}(t,w,d,h)=(\Uc(d,h),w,0,t).
\]
The following words implement the three stack operations.
\[
\begin{aligned}
\text{pop}\quad & \ax(0,0,0)\,\shift^{10},\\
\text{push}\quad & \ax(0,0,0)\,\shift^{8}\,\mpc(0),\\
\text{discard}\quad & \mpc(1)\,\ax(1,0,0)\,\shift^{9}\,\mpc(0).
\end{aligned}
\]
Each formula these Shifts test is $0$ or at least $2$. The heads
and gate have residue $0$, and $\Uc(d,h)$ has residue $2$. To select
head $j$ from the top, for $1\le j\le7$, discard, call $\ax$, then
apply $j+1$ Shifts and $8-j$ pops. To select the gate, discard, call
$\ax$, then pop nine times. Discard sets $d=h=0$, so
$(D,H,A)=(x,y,c)$. Read the corresponding row of
Table~\ref{tab:heads} backwards to choose the arguments for an
instance $\rho(\alpha,\beta,\gamma)$. For $\rho_8(a,b,c,u)$ with
$u\le a$, set the offset to $u$ and call $\ax(b,a-u,c)$.

The generalization word pushes, applies offset $i$ and a Shift,
calls $\ax(i,0,0)$ and applies $\shift^{9}$. The modus ponens suffix
is $\shift^{2}\,\mpc(b)\,\ax(b,0,0)\,\shift^{10}$. With these words,
the induction on fuel and the scheduler argument
of~\cite[\S6.2--6.3]{pa209} apply as written.
\end{proof}

\section{The layout}\label{sec:layout}

At a node, the state reads the path bit, moves right and enters the
state for the corresponding child. Leaves are instructions, and
the machine stores equal subtrees once. The 201 items in the body of
$\ax$ form the leaves of a binary tree with 68 distinct internal
nodes.

As in $M_{209}$, two nodes can share a state if one occurs only as a
left child and the other only as a right child. When the current bit
would select different children for the two nodes, the shared state
moves left to the bit that selected the node. A second state reads
that bit again and moves right into a third state, which uses the
current bit to choose the correct child. Both states already read
other nodes or start a macro, so sharing requires no extra states.
The six pairs in Table~\ref{tab:shared} let 62 states read the 68
nodes.

\begin{table}[ht]
\centering\small
\begin{tabular}{@{}clll@{}}
\toprule
State & Left-child node & Right-child node & Rows sent left \\
\midrule
138 & $(\mathtt{D}\mid\mathtt{HC})$ & $(\mathtt{PiCDCiiPiC}\mid\mathtt{HC})$ & $0\to168$ \\
140 & $(\mathtt{C}\mid\mathtt{C})$ & $(\mathtt{H}\mid\mathtt{C})$ & $0\to158$ \\
160 & $(\mathtt{i}\mid\mathtt{ADCi})$ & $(\mathtt{i}\mid\mathtt{C})$ & $1\to185$ \\
174 & $(\mathtt{i}\mid\mathtt{ACiiCHDCi})$ & $(\mathtt{DCiiHDCii}\mid\mathtt{ACiiCHDCi})$ & $0\to172$ \\
179 & $(\mathtt{A}\mid\mathtt{ADCi})$ & $(\mathtt{CC}\mid\mathtt{ADCi})$ & $0\to142$ \\
183 & $(\mathtt{C}\mid\mathtt{PiC})$ & $(\mathtt{H}\mid\mathtt{DCi})$ & $0\to158$, $1\to167$ \\
\bottomrule
\end{tabular}
\caption{States of $M$ that read two nodes of the body of $\ax$.
Each node lists its two children as strings of items. \texttt{P}
pushes. \texttt{D}, \texttt{H} and \texttt{A} push and load a
register. \texttt{C} pops and forms $4p$, and \texttt{i} adds one.
The last column gives each symbol whose row moves left and the
state it enters.}\label{tab:shared}
\end{table}

We remove two states outside the body that $M_{209}$ used to
distinguish nodes of its old body. Each read a single symbol and
stepped back onto the cell just written. The row that entered it
now performs the successor's action directly.

\section{Simulation}\label{sec:simulation}

Number the cells of a frame from the first path cell. In a
\emph{clean} frame, zeros separate the path from the header $1$, and
the register file follows the header. The \emph{gap} is the position
of that header. The longest path in $R$ is 38 bits. The formal
simulation uses reachable frames with a gap of at least 48 and a
file of at least ten registers. The clean startup frame meets the
gap bound, and the simulation preserves it.

\begin{lemma}[One instruction]\label{lem:contract}
For every site of $R$ and each outcome of a conditional decrement,
the machine $M$ started in state $75$ on a valid clean frame of
that site with gap at least 48 reaches, in at least one step, the clean frame of the
successor site with the register effect of the instruction. At the
halt site it enters the halt label.
\end{lemma}

\begin{proof}
A run from a clean frame alternates between walks outside the engine
and calls to the engine. The engine rows match those of $M_{209}$.
The register lemmas of~\cite[\S3]{pa209} apply for the same register
contents and gaps. They cover filling the gap, marking a register,
incrementing, both outcomes of decrementing, updating the path and
scanning back. The return scan ends on the cell after the first $1$
with seven zeros to its left. For most sites that cell is the root.
At 19 sites one pass ends on a later path cell, where the next walk
begins, as in $M_{209}$.

In the computed runs, states outside the engine read only path
cells and the two cells after the path. These reads stay at or to the
right of the root. At each site and for each outcome, the walk visits
the same sequence of states, cells and symbols for every register
file tested. One run thus determines the fixed sequence of rows
and engine calls for the contract.

We ran the checks at all $3{,}812$ sites, including both outcomes
at each of the $1{,}147$ decrement sites. We varied the register
contents and used gaps $40$, $41$, $48$ and $61$. All $68{,}616$
runs end in the required frame, and together they read all 392
rows.
\end{proof}

The coverage theorem in \file{PA196/Coverage.lean} checks the
positive-length contract at every site for the reachable gap
bound, including both outcomes of a conditional decrement.

\begin{lemma}[Startup]\label{lem:startup}
From the blank tape $M$ reaches after exactly $58{,}154$ steps the
clean frame of site $25$ with gap $48$, candidate register $2$, every
other logical register $0$, and $26$ further registers equal to $1$.
\end{lemma}

\begin{proof}
After running $M$ for $58{,}154$ steps, the tape matches the startup
frame of $M_{209}$ cell for cell~\cite[\S2]{pa209}. The head is on
the first path cell in state $75$. Lean checks each segment and
joins them to obtain the literal startup trace.
\end{proof}

The scheduler starts at site $25$ with candidate $2$ and every
other logical register zero, the clean store of~\cite{pa209}.
\file{PA196/Startup.lean} verifies the startup, and
\file{PA196/Nonhalting.lean} connects it to that store.

\begin{proof}[Proof of Theorem~\ref{thm:main}]
Lemmas~\ref{lem:contract} and~\ref{lem:startup} identify halting of
$M$ from the blank tape with halting of $R$ from the clean store.
Under PA consistency, Lemma~\ref{lem:sound} gives nonhalting of $R$,
so $M$ runs forever and $\neg H$ is true in $\N$.

In any model of $\PA$, the compiler turns a coded PA refutation into
a halting command word. The scheduler reaches that word from
candidate $2$. The tape simulation joins the startup to the
instruction traces, then projects the internal control codes onto
the literal rows. This trace witnesses $H$, including when the proof
code and trace length are nonstandard. Hence
\[
   \PA\vdash \Prov_{\PA}(\ulcorner\bot\urcorner)\to H,
   \qquad
   \PA\vdash \neg H\to\Con(\PA).
\]
G\"odel's second incompleteness theorem~\cite{goedel} excludes a
$\PA$ proof of $\neg H$. If $\PA$ is sound, it does not prove the
false sentence $H$.
\end{proof}

\clearpage
\section{Lean verification}\label{sec:status}

Lean checks all 392 transitions, the startup trace and the
simulation contracts at all 3,812 source sites.
The formal contracts use the reachable lower bound 48, which
suffices for the blank-tape theorem. The finite runs at gaps 40 and
41 give additional checks.

\file{PA196/Arithmetic/RefutationHalting.lean} connects the arithmetic
trace construction to halting. The trace sentence uses the 196
working states and the distinct halt label 196.
\file{PA196/Independence.lean} proves the internal consistency
implication and two-sided independence. \file{PA196/Nonhalting.lean}
proves the machine theorem.

\file{Challenge.lean} states nonhalting and independence using an
explicit PA proof calculus and the literal table.
\file{PA196/Palomar/Solution.lean} proves those declarations and
connects their definitions to Foundation and the tape model. Guarded
axiom audits cover both presentations. The repository includes the
table and the checked proof sources. \file{docs/reproduction.md}
gives the finite checks, Lean build and Comparator commands.

\begin{thebibliography}{9}

\bibitem{pa209}
J. Chang.
Nonhalting and PA-independence for a 209-state Turing machine.
\href{https://github.com/jamesyc/BusyBeaver/tree/69d98241a721ead20b4f0b17df9062950b40fbfe}{Public PA209 snapshot}.

\bibitem{goedel}
K. G\"odel.
\"Uber formal unentscheidbare S\"atze der Principia Mathematica und
verwandter Systeme I.
\emph{Monatshefte f\"ur Mathematik und Physik} 38 (1931), 173--198.
\url{https://doi.org/10.1007/BF01700692}.

\end{thebibliography}
\end{document}
